\section{Local3D URANS--VOF Model}

Local3D hydrodynamics are implemented through the adopted OpenFOAM Foundation 11 incompressible two-phase URANS--VOF backend. The repository owns the case-generation policy and validation wrapper; it does not reimplement OpenFOAM solver internals.

The local continuity equation is

\begin{equation}
\label{eq:local-continuity}
\nabla\cdot\mathbf u=0.
\end{equation}

The two-phase URANS momentum equation is represented as

\begin{equation}
\label{eq:urans-momentum}
\frac{\partial \rho\mathbf u}{\partial t}
+
\nabla\cdot(\rho\mathbf u\otimes\mathbf u)
=
-
\nabla p
+
\nabla\cdot\left[
(\mu+\mu_t)
\left(\nabla\mathbf u+\nabla\mathbf u^{T}\right)
\right]
+
\rho\mathbf g .
\end{equation}

The VOF phase-fraction equation is

\begin{equation}
\label{eq:vof-transport}
\frac{\partial\alpha}{\partial t}
+
\nabla\cdot(\alpha\mathbf u)
+
\nabla\cdot
\left[
\alpha(1-\alpha)\mathbf u_c
\right]
=
0 .
\end{equation}

Mixture properties are

\begin{equation}
\label{eq:mixture-density}
\rho(\alpha)=\alpha\rho_w+(1-\alpha)\rho_a
\end{equation}

and

\begin{equation}
\label{eq:mixture-viscosity}
\mu(\alpha)=\alpha\mu_w+(1-\alpha)\mu_a .
\end{equation}

The SST turbulence model is represented by transport equations for \(k\) and \(\omega\), with the adopted backend supplying the detailed coefficients \autocite{menter1994twoequation}:

\begin{equation}
\label{eq:sst-k}
\frac{\partial k}{\partial t}
+
\mathbf u\cdot\nabla k
=
P_k-\beta^\ast k\omega
+
\nabla\cdot\left[(\nu+\sigma_k\nu_t)\nabla k\right]
\end{equation}

and

\begin{equation}
\label{eq:sst-omega}
\frac{\partial \omega}{\partial t}
+
\mathbf u\cdot\nabla\omega
=
\gamma\frac{\omega}{k}P_k-\beta\omega^2
+
\nabla\cdot\left[(\nu+\sigma_\omega\nu_t)\nabla\omega\right]
+D_\omega .
\end{equation}

The eddy viscosity closure is

\begin{equation}
\label{eq:eddy-viscosity}
\nu_t=\frac{a_1k}{\max(a_1\omega,SF_2)} .
\end{equation}

Barrier traction combines pressure and viscous stress:

\begin{equation}
\label{eq:barrier-traction}
\mathbf t=
\boldsymbol\sigma\mathbf n_b
=
\left[-p\mathbf I+\boldsymbol\tau_v\right]\mathbf n_b .
\end{equation}

The force is

\begin{equation}
\label{eq:barrier-force}
\mathbf F
=
\int_{\Gamma_b}
\mathbf t\,\mathrm dA
\end{equation}

and the moment is

\begin{equation}
\label{eq:barrier-moment}
\mathbf M
=
\int_{\Gamma_b}
(\mathbf x-\mathbf x_0)
\times
\mathbf t
\,\mathrm dA .
\end{equation}

The future impulse metric is

\begin{equation}
\label{eq:future-impulse}
\mathbf J
=
\int_{t_1}^{t_2}
\mathbf F(t)\,\mathrm dt .
\end{equation}

Local3D turbulence-model sensitivity is methodology-relevant \autocite{kozelkov2022turbulence,watanabe2022validation}, but physical calibration has not started.
